\chapter{Conclusion}

This work implemented and tested fourteen algorithm configurations for the
Linear Ordering Problem --- twelve iterative improvement variants and two VND
algorithms --- on 78 benchmark instances of sizes 150 and 250.

The most clear result is that the neighbourhood structure matters more than
anything else. Insert-based search consistently reaches around 2\% average
RPD, exchange lands somewhere between 2.8\% and 3.9\%, and transpose is
essentially unusable on its own, reaching 17--35\% depending on
initialisation. These differences are all statistically significant and
persistent across instances, so this is not a fluke of the benchmark set.

The pivoting rule also has a noticeable effect, but in a direction that is not
immediately obvious: first-improvement beats best-improvement for both insert
and exchange. The most likely explanation is that best-improvement makes
fewer, greedier moves and gets trapped in worse local optima, while
first-improvement explores the search space more gradually. For transpose, the
pivoting rule does not matter much since the neighbourhood is too small either
way.

Initialisation is only relevant when using transpose or exchange. For insert,
starting from a random permutation or from the CW heuristic gives essentially
the same final result ($p > 0.5$). The insert neighbourhood is large enough to
recover from a poor starting point, which is a useful practical observation.

The VND results are perhaps the least expected. Combining all three
neighbourhoods does not significantly outperform FI-insert-CW on its own
($p = 0.636$ for VND1, $p = 0.053$ for VND2). The insert neighbourhood
already covers most of what the others bring, so the added complexity of VND
gives limited benefit here. The two VND orderings are also statistically
indistinguishable from each other ($p = 0.176$).

If we had to pick one algorithm to use in practice, FI-insert-CW is a solid
choice: it matches VND in quality, runs in a reasonable time, and is simpler
to implement. BI-insert is an option if speed is more important, though it
gives slightly worse results.

One limitation worth mentioning is that each configuration is run only once
per instance. For the random initialisation variants, multiple runs would give
a more reliable estimate of average performance. A natural follow-up would
also be to wrap these local search algorithms inside an iterated local search
or simulated annealing framework to escape local optima and push the RPD
further down.
